% libre-compendium
%
% Copyright (C) 2026 Richard Qin
%
% This document is dual-licensed under either:
% - GNU Free Documentation License, Version 1.3 or any later version
%   published by the Free Software Foundation; with no Invariant Sections,
%   no Front-Cover Texts, and no Back-Cover Texts (see COPYING.DOC); or
% - Creative Commons Attribution-ShareAlike 4.0 International License
%   (see LICENSE-CC-BY-SA-4.0).

\chapter{二次型、对称双线性型与正定矩阵}

\section{研究范围与预备事实}

本章主要在实数域 $\R$ 上讨论二次型、实对称矩阵和对称双线性型。若把结论推广到一般数域 $\F$，至少应假设 $\operatorname{char}\F\ne 2$，因为从交叉项系数恢复对称矩阵以及极化公式都需要除以 $2$。正定、负定、正惯性指数等带有大小关系的概念只在有序域上有意义，本章统一放在 $\R$ 上陈述。

对矩阵 $A=(a_{ij})\in\R^{n\times n}$，若 $A^{\mathsf T}=A$，则称 $A$ 为实对称矩阵。实对称矩阵具有两个基本性质：其特征值均为实数；存在正交矩阵 $U$ 使
\[
 U^{\mathsf T}AU=\diag(\lambda_1,\ldots,\lambda_n).
\]
这里 $U^{-1}=U^{\mathsf T}$，$\lambda_i$ 为 $A$ 的全部特征值，按代数重数计数。

\section{二次型及其对称矩阵}

\begin{definition}[二次型]
设 $A\in\R^{n\times n}$ 为实对称矩阵，$x=(x_1,\ldots,x_n)^{\mathsf T}$。函数
\[
 Q:\R^n\to\R,\qquad Q(x)=x^{\mathsf T}Ax
\]
称为 $n$ 元实二次型，$A$ 称为 $Q$ 在标准基下的矩阵。
\end{definition}

展开矩阵乘法可得
\[
Q(x)=\sum_{i=1}^n a_{ii}x_i^2+2\sum_{1\le i<j\le n}a_{ij}x_ix_j.
\]
因此，平方项 $x_i^2$ 的系数就是 $a_{ii}$，交叉项 $x_ix_j$ 的系数是 $2a_{ij}$。由多项式写矩阵时，交叉项系数必须折半后分别放入 $(i,j)$ 与 $(j,i)$ 位置。

\begin{proposition}[二次齐次多项式与对称矩阵的一一对应]
每个实二次齐次多项式
\[
 Q(x)=\sum_{i=1}^n b_{ii}x_i^2+\sum_{1\le i<j\le n}b_{ij}x_ix_j
\]
唯一对应一个实对称矩阵 $A$，其中
\[
 a_{ii}=b_{ii},\qquad a_{ij}=a_{ji}=\frac{b_{ij}}2\quad(i<j).
\]
\end{proposition}

\begin{proof}
上述取法直接给出 $Q(x)=x^{\mathsf T}Ax$，故存在性成立。若另有对称矩阵 $C$ 也满足 $Q(x)=x^{\mathsf T}Cx$，比较各平方项和交叉项的系数，得到 $c_{ii}=a_{ii}$ 以及 $2c_{ij}=2a_{ij}$，故 $C=A$。
\end{proof}

\begin{example}[title={由多项式写出矩阵}]
求二次型
\[
 Q(x_1,x_2,x_3)=2x_1^2-6x_1x_2+4x_1x_3+5x_2^2-2x_2x_3+x_3^2
\]
对应的对称矩阵。
\begin{solution}
交叉项系数折半，得
\[
A=\begin{pmatrix}
2&-3&2\\
-3&5&-1\\
2&-1&1
\end{pmatrix}.
\]
例如 $x_1x_2$ 项由 $a_{12}x_1x_2+a_{21}x_2x_1=2a_{12}x_1x_2$ 产生，因此 $a_{12}=-3$。
\end{solution}
\end{example}

\begin{pitfall}[交叉项不折半]
把 $b_{ij}$ 直接填入 $a_{ij}$ 会使展开后的交叉项变成 $2b_{ij}x_ix_j$。二次型矩阵默认取对称矩阵；非对称矩阵虽然也可写出同一个 $x^{\mathsf T}Ax$，但只有其对称部分真正起作用。
\end{pitfall}

事实上，对任意方阵 $M$，令
\[
S=\frac{M+M^{\mathsf T}}2,\qquad K=\frac{M-M^{\mathsf T}}2,
\]
则 $S$ 对称、$K$ 反对称，并且 $x^{\mathsf T}Kx=0$。故
\[
x^{\mathsf T}Mx=x^{\mathsf T}Sx.
\]
这说明要求二次型矩阵对称并非额外限制，而是为了得到唯一表示。

\section{二次型与对称双线性型}

\begin{definition}[双线性型与对称双线性型]
设 $V$ 是 $\R$ 上的线性空间。映射 $B:V\times V\to\R$ 若分别对两个变元线性，即对任意 $u,v,w\in V$ 和 $a,b\in\R$ 满足
\[
 B(au+bv,w)=aB(u,w)+bB(v,w),
\]
\[
 B(w,au+bv)=aB(w,u)+bB(w,v),
\]
则称 $B$ 为双线性型。若进一步满足 $B(u,v)=B(v,u)$，则称 $B$ 为对称双线性型。
\end{definition}

给定实对称矩阵 $A$，定义
\[
B(x,y)=x^{\mathsf T}Ay.
\]
矩阵乘法的分配律与数乘相容性保证 $B$ 对两个变元分别线性；又因为 $A^{\mathsf T}=A$，
\[
B(y,x)=y^{\mathsf T}Ax=(y^{\mathsf T}Ax)^{\mathsf T}=x^{\mathsf T}A^{\mathsf T}y=x^{\mathsf T}Ay=B(x,y).
\]
其对角限制恰为二次型：$Q(x)=B(x,x)$。

\begin{theorem}[实数域上的极化公式]
若 $Q(x)=B(x,x)$，其中 $B$ 为对称双线性型，则
\[
B(x,y)=\frac{Q(x+y)-Q(x)-Q(y)}2
=\frac{Q(x+y)-Q(x-y)}4.
\]
因此，一个实二次型唯一确定一个对称双线性型。
\end{theorem}
\begin{proof}
由双线性和对称性，
\[
Q(x+y)=B(x+y,x+y)=Q(x)+2B(x,y)+Q(y),
\]
整理得到第一式。类似地，
\[
Q(x-y)=Q(x)-2B(x,y)+Q(y),
\]
两式相减得到第二式。
\end{proof}

\begin{remark}
表达式 $[Q(x+y)-Q(y)]/2$ 通常不等于 $B(x,y)$，因为它还包含 $Q(x)/2$。只有把 $Q(x)$ 一并减去，才是正确的极化公式。
\end{remark}

\begin{example}[title={由二次型恢复双线性型}]
设 $Q(x_1,x_2)=3x_1^2-4x_1x_2+2x_2^2$。求对应的 $B(x,y)$。
\begin{solution}
二次型矩阵为
\[
A=\begin{pmatrix}3&-2\\-2&2\end{pmatrix}.
\]
故对 $x=(x_1,x_2)^{\mathsf T}$、$y=(y_1,y_2)^{\mathsf T}$，
\[
B(x,y)=x^{\mathsf T}Ay=3x_1y_1-2x_1y_2-2x_2y_1+2x_2y_2.
\]
将 $y=x$ 代入即可检验 $B(x,x)=Q(x)$。
\end{solution}
\end{example}

\section{双线性型的矩阵与基变换}

设 $V$ 为 $n$ 维实线性空间，$\mathcal E=(e_1,\ldots,e_n)$ 是一组有序基。对双线性型 $B$，定义
\[
G_{\mathcal E}=(B(e_i,e_j))_{n\times n}.
\]
该矩阵称为 $B$ 在基 $\mathcal E$ 下的矩阵，也常称为 Gram 矩阵。这里的“Gram 矩阵”取决于所讨论的双线性型和基；它不必是欧氏内积的 Gram 矩阵。

\begin{theorem}[有限维双线性型的矩阵表示]
给定 $n$ 维实线性空间 $V$ 的一组有序基 $\mathcal E$，每个双线性型 $B:V\times V\to\R$ 都唯一对应一个矩阵 $G\in\R^{n\times n}$，使得对任意 $u,v\in V$，
\[
B(u,v)=[u]_{\mathcal E}^{\mathsf T}G[v]_{\mathcal E}.
\]
矩阵元为 $g_{ij}=B(e_i,e_j)$。此外，$B$ 对称当且仅当 $G$ 对称。
\end{theorem}
\begin{proof}
先证存在性。设
\[
u=\sum_{i=1}^n x_ie_i,\qquad v=\sum_{j=1}^n y_je_j.
\]
固定第二个变元并利用第一变元的线性性，再固定第一变元并利用第二变元的线性性，得到
\begin{align*}
B(u,v)
&=B\left(\sum_{i=1}^n x_ie_i,\sum_{j=1}^n y_je_j\right)\\
&=\sum_{i=1}^n x_iB\left(e_i,\sum_{j=1}^n y_je_j\right)\\
&=\sum_{i=1}^n\sum_{j=1}^n x_iy_jB(e_i,e_j).
\end{align*}
令 $G=(g_{ij})$，其中 $g_{ij}=B(e_i,e_j)$，上式正是 $x^{\mathsf T}Gy$。

再证唯一性。若矩阵 $H$ 也满足 $B(u,v)=x^{\mathsf T}Hy$，分别取 $u=e_i$、$v=e_j$。此时坐标向量为第 $i,j$ 个标准单位向量，故
\[
h_{ij}=e_i^{\mathsf T}He_j=B(e_i,e_j)=g_{ij}.
\]
对所有 $i,j$ 都成立，所以 $H=G$。

若 $B$ 对称，则 $g_{ij}=B(e_i,e_j)=B(e_j,e_i)=g_{ji}$，故 $G^{\mathsf T}=G$。反之，若 $G$ 对称，则
\[
B(v,u)=y^{\mathsf T}Gx=(y^{\mathsf T}Gx)^{\mathsf T}
=x^{\mathsf T}G^{\mathsf T}y=x^{\mathsf T}Gy=B(u,v),
\]
故 $B$ 对称。
\end{proof}

因此，若 $[u]_{\mathcal E}=x$、$[v]_{\mathcal E}=y$，则
\[
B(u,v)=\sum_{i,j=1}^n x_i y_j B(e_i,e_j)=x^{\mathsf T}G_{\mathcal E}y.
\]
若 $B$ 对称，则 $G_{\mathcal E}$ 对称；反之亦然。

\begin{theorem}[基变换下的矩阵变换]
设新基 $\widetilde{\mathcal E}=(\widetilde e_1,\ldots,\widetilde e_n)$ 与旧基 $\mathcal E$ 满足
\[
(\widetilde e_1,\ldots,\widetilde e_n)=(e_1,\ldots,e_n)P,
\]
其中 $P$ 可逆。则
\[
[u]_{\mathcal E}=P[u]_{\widetilde{\mathcal E}},\qquad
G_{\widetilde{\mathcal E}}=P^{\mathsf T}G_{\mathcal E}P.
\]
\end{theorem}
\begin{proof}
矩阵等式 $\widetilde{\mathcal E}=\mathcal E P$ 的含义是：$P$ 的第 $j$ 列等于新基向量 $\widetilde e_j$ 在旧基 $\mathcal E$ 下的坐标。记 $\widetilde x=[u]_{\widetilde{\mathcal E}}$，则
\[
u=\widetilde{\mathcal E}\widetilde x=\mathcal E P\widetilde x,
\]
另一方面 $u=\mathcal E x$。由于向量在同一组基下的坐标唯一，故旧坐标 $x=P\widetilde x$。同理 $y=P\widetilde y$。于是
\[
B(u,v)=x^{\mathsf T}G_{\mathcal E}y
=\widetilde x^{\mathsf T}(P^{\mathsf T}G_{\mathcal E}P)\widetilde y.
\]
若新基下的矩阵记为 $\widetilde G$，则同一个数 $B(u,v)$ 还等于 $\widetilde x^{\mathsf T}\widetilde G\widetilde y$。因此对任意 $\widetilde x,\widetilde y\in\R^n$，
\[
\widetilde x^{\mathsf T}\bigl(\widetilde G-P^{\mathsf T}G_{\mathcal E}P\bigr)\widetilde y=0.
\]
依次取 $\widetilde x,\widetilde y$ 为标准单位向量，可知括号内矩阵的每个元素都为零，故 $\widetilde G=P^{\mathsf T}G_{\mathcal E}P$。
\end{proof}

\begin{pitfall}[基变换与坐标变换的方向]
矩阵 $P$ 的第 $j$ 列是新基向量 $\widetilde e_j$ 在旧基下的坐标。因此“新基等于旧基乘 $P$”时，旧坐标等于 $P$ 乘新坐标：$x=P\widetilde x$；反向则是 $\widetilde x=P^{-1}x$。
\end{pitfall}

\begin{example}[title={计算新基下的 Gram 矩阵}]
在 $\R^2$ 上，设
\[
B(x,y)=2x_1y_1+3x_1y_2+3x_2y_1+5x_2y_2.
\]
取新基 $\widetilde e_1=(1,1)^{\mathsf T}$、$\widetilde e_2=(1,-1)^{\mathsf T}$，求新基下的矩阵。
\begin{solution}
标准基下
\[
G=\begin{pmatrix}2&3\\3&5\end{pmatrix},\qquad
P=\begin{pmatrix}1&1\\1&-1\end{pmatrix}.
\]
故
\[
\widetilde G=P^{\mathsf T}GP
=\begin{pmatrix}13&-3\\-3&1\end{pmatrix}.
\]
也可直接计算 $B(\widetilde e_i,\widetilde e_j)$ 验证四个矩阵元。
\end{solution}
\end{example}

\section{矩阵相合与变量替换}

\begin{definition}[矩阵相合]
设 $A,B\in\R^{n\times n}$。若存在可逆矩阵 $P$ 使
\[
B=P^{\mathsf T}AP,
\]
则称 $A$ 与 $B$ 相合，记作 $A\cge B$。
\end{definition}

相合关系是等价关系：自反性取 $P=I$；对称性由
\[
A=(P^{-1})^{\mathsf T}BP^{-1}
\]
得到；若 $B=P^{\mathsf T}AP$ 且 $C=Q^{\mathsf T}BQ$，则
\[
C=(PQ)^{\mathsf T}A(PQ),
\]
故具有传递性。

\begin{proposition}[相合变换的基础不变量]
若 $B=P^{\mathsf T}AP$，其中 $P$ 可逆，则
\[
\rank B=\rank A,\qquad
\det B=(\det P)^2\det A.
\]
因此，相合保持秩、可逆性以及非零行列式的符号，但通常不保持行列式的具体数值和特征值。
\end{proposition}
\begin{proof}
矩阵 $P^{\mathsf T}$ 与 $P$ 都可逆，左右乘可逆矩阵不改变秩，故
\[
\rank(P^{\mathsf T}AP)=\rank A.
\]
利用行列式的乘法公式和 $\det P^{\mathsf T}=\det P$，得到
\[
\det B=\det(P^{\mathsf T})\det A\det P
=(\det P)^2\det A.
\]
由于 $(\det P)^2>0$，当 $\det A\ne0$ 时，$\det A$ 与 $\det B$ 同号。
\end{proof}

对二次型 $Q(x)=x^{\mathsf T}Ax$ 作可逆变量替换 $x=Py$，则
\[
Q(x)=(Py)^{\mathsf T}A(Py)=y^{\mathsf T}(P^{\mathsf T}AP)y.
\]
因此，变量替换后的二次型矩阵与原矩阵相合。相合是“同一双线性型或二次型在不同基、不同坐标下”的矩阵语言。

\begin{summary}[相合、相似与等价的区别]
相合变换为 $P^{\mathsf T}AP$，研究二次型和双线性型；相似变换为 $P^{-1}AP$，研究同一线性变换在不同基下的矩阵；矩阵等价变换为 $PAQ$，研究线性映射在定义域、值域分别换基后的矩阵。相似保持特征值，相合一般不保持特征值，但保持秩及惯性指数。
\end{summary}

\begin{example}[title={相合不保持特征值}]
取 $A=(1)$、$P=(2)$，则 $P^{\mathsf T}AP=(4)$。矩阵 $(1)$ 与 $(4)$ 相合，但特征值分别为 $1$ 与 $4$。
\end{example}

\section{配方法与相合对角化}

\subsection{首个平方项系数非零的情形}

设
\[
Q(x)=\sum_{i,j=1}^n a_{ij}x_ix_j,\qquad A=(a_{ij})=A^{\mathsf T},
\]
且 $a_{11}\ne0$。集中所有含 $x_1$ 的项：
\begin{align*}
Q(x)
&=a_{11}x_1^2+2x_1\sum_{j=2}^n a_{1j}x_j+Q_0(x_2,\ldots,x_n)\\
&=a_{11}\left(x_1+\sum_{j=2}^n\frac{a_{1j}}{a_{11}}x_j\right)^2
+Q_1(x_2,\ldots,x_n).
\end{align*}
令
\[
y_1=x_1+\sum_{j=2}^n\frac{a_{1j}}{a_{11}}x_j,\qquad y_i=x_i\ (i\ge2),
\]
得到 $Q=a_{11}y_1^2+Q_1(y_2,\ldots,y_n)$。该替换的系数矩阵为三角矩阵且对角元全为 $1$，故可逆。

在矩阵语言中，若
\[
A=\begin{pmatrix}a&\alpha^{\mathsf T}\\\alpha&A_1\end{pmatrix},\qquad a\ne0,
\]
取
\[
P=\begin{pmatrix}1&-\alpha^{\mathsf T}/a\\0&I_{n-1}\end{pmatrix},
\]
则
\[
P^{\mathsf T}AP=
\begin{pmatrix}
a&0\\
0&A_1-\dfrac{\alpha\alpha^{\mathsf T}}a
\end{pmatrix}.
\]
右下角矩阵称为关于主元 $a$ 的 Schur 补。

\subsection{首个平方项系数为零的处理}

若 $a_{11}=0$，不能除以 $a_{11}$。有两类处理：
\begin{enumerate}
\item 若存在 $a_{ii}\ne0$，交换 $x_1$ 与 $x_i$，把非零平方项移到首位。
\item 若所有 $a_{ii}=0$ 但 $A\ne0$，则存在 $a_{ij}\ne0$。交换变量后可设 $a_{12}\ne0$，再令
\[
x_1=y_1+y_2,\qquad x_2=y_1-y_2.
\]
此时 $2a_{12}x_1x_2=2a_{12}(y_1^2-y_2^2)$，产生非零平方项，转化为前一种情形。
\end{enumerate}

\begin{theorem}[对称矩阵的相合对角化]
任意实对称矩阵 $A$ 都相合于某个实对角矩阵。等价地，任意实二次型都可经可逆线性变量替换化为平方项之和。
\end{theorem}
\begin{proof}
对矩阵阶数 $n$ 作归纳。

当 $n=1$ 时，任意矩阵 $A=(a_{11})$ 已经是对角矩阵。

假设任意 $(n-1)$ 阶实对称矩阵都能相合对角化，考虑 $n$ 阶实对称矩阵 $A$。若 $A=0$，结论成立。以下设 $A\ne0$，分三步处理。

\textbf{第一步：取得非零左上角元素。} 若 $a_{11}\ne0$，无需变换。若 $a_{11}=0$ 但某个 $a_{ii}\ne0$，用置换矩阵同时交换第 $1$ 行与第 $i$ 行、第 $1$ 列与第 $i$ 列；这正是相合变换 $S^{\mathsf T}AS$，且新矩阵左上角为 $a_{ii}\ne0$。若所有对角元均为零，由 $A\ne0$ 可知存在 $i\ne j$ 使 $a_{ij}\ne0$。先用置换把它移到 $(1,2)$ 位置，再把第一、第二个变量替换为
\[
x_1=y_1+y_2,\qquad x_2=y_1-y_2.
\]
因为原来的 $2a_{12}x_1x_2$ 变为 $2a_{12}(y_1^2-y_2^2)$，变换后二次型的 $y_1^2$ 系数非零。该变量替换的系数矩阵行列式为 $-2$，故可逆。

因此，经过某个可逆相合变换后，总可设矩阵具有分块形式
\[
A=\begin{pmatrix}a&\alpha^{\mathsf T}\\\alpha&A_1\end{pmatrix},
\qquad a\ne0.
\]

\textbf{第二步：消去第一行和第一列的非对角元。} 取
\[
P_1=\begin{pmatrix}1&-\alpha^{\mathsf T}/a\\0&I_{n-1}\end{pmatrix}.
\]
$P_1$ 是对角元均为 $1$ 的分块上三角矩阵，所以 $\det P_1=1$，从而可逆。直接分块相乘：
\[
AP_1=
\begin{pmatrix}
a&0\\
\alpha&A_1-\alpha\alpha^{\mathsf T}/a
\end{pmatrix},
\]
继而
\[
P_1^{\mathsf T}AP_1=
\begin{pmatrix}
a&0\\
0&A_1-\alpha\alpha^{\mathsf T}/a
\end{pmatrix}.
\]
右下角矩阵仍然对称，因为 $A_1$ 与 $\alpha\alpha^{\mathsf T}$ 都对称。

\textbf{第三步：对低一阶矩阵使用归纳假设。} 记
\[
A_1'=A_1-\frac{\alpha\alpha^{\mathsf T}}a.
\]
由归纳假设，存在可逆矩阵 $R$，使 $R^{\mathsf T}A_1'R=D$ 为对角矩阵。令
\[
P_2=\diag(1,R),
\]
则 $P_2$ 可逆，且
\[
P_2^{\mathsf T}
\begin{pmatrix}a&0\\0&A_1'\end{pmatrix}
P_2
=\begin{pmatrix}a&0\\0&D\end{pmatrix},
\]
右端为对角矩阵。把第一步的变换、$P_1$ 与 $P_2$ 依次相乘，仍得到可逆矩阵，因此原矩阵 $A$ 相合于对角矩阵。
\end{proof}

\begin{corollary}[秩与非零平方项个数]
若二次型经可逆变量替换化为 $d_1y_1^2+\cdots+d_ry_r^2$，其中每个 $d_i\ne0$，其余变量的系数为零，则 $r=\rank A$。
\end{corollary}
\begin{proof}
规范形矩阵为 $D=\diag(d_1,\ldots,d_r,0,\ldots,0)$，其秩为 $r$。相合变换是在矩阵左右分别乘可逆矩阵 $P^{\mathsf T}$ 与 $P$，不改变秩，故 $\rank A=\rank D=r$。
\end{proof}

\begin{definition}[二次型的秩]
二次型 $Q(x)=x^{\mathsf T}Ax$ 的秩定义为其对称矩阵 $A$ 的秩。由上一个推论，这也等于任意对角形中非零平方项的个数，因此不依赖坐标或配方方法的选择。
\end{definition}

\begin{example}[title={三元二次型的逐步配方}]
化简
\[
G=x_1^2+2x_2^2+3x_3^2+4x_1x_2+2x_1x_3+2x_2x_3.
\]
\begin{solution}
先对 $x_1$ 配方：
\begin{align*}
G&=(x_1+2x_2+x_3)^2-2x_2^2-2x_2x_3+2x_3^2\\
&=(x_1+2x_2+x_3)^2-2\left(x_2+\frac12x_3\right)^2+\frac52x_3^2.
\end{align*}
令
\[
z_1=x_1+2x_2+x_3,\quad z_2=x_2+\frac12x_3,\quad z_3=x_3,
\]
则
\[
G=z_1^2-2z_2^2+\frac52z_3^2.
\]
这里写出的是 $z=Cx$。若要求相合矩阵 $P$ 满足 $x=Pz$，则 $P=C^{-1}$，不能直接把 $C$ 当作 $P$。
\end{solution}
\end{example}

\begin{example}[title={没有平方项时的配方}]
将 $Q(x_1,x_2)=2x_1x_2$ 化为对角形。
\begin{solution}
令 $y_1=x_1+x_2$、$y_2=x_1-x_2$，则
\[
x_1=\frac{y_1+y_2}{2},\qquad x_2=\frac{y_1-y_2}{2},
\]
所以
\[
Q=\frac12y_1^2-\frac12y_2^2.
\]
继续缩放变量 $u_1=y_1/\sqrt2$、$u_2=y_2/\sqrt2$，可得规范形 $Q=u_1^2-u_2^2$。
\end{solution}
\end{example}

\section{正交对角化与相合对角化}

若 $A$ 为实对称矩阵，谱定理给出正交矩阵 $U$，使
\[
U^{\mathsf T}AU=\diag(\lambda_1,\ldots,\lambda_n).
\]
变量替换 $x=Uy$ 保持欧氏长度，因为 $\lVert x\rVert^2=y^{\mathsf T}U^{\mathsf T}Uy=\lVert y\rVert^2$。此时对角元必须是特征值。

一般相合对角化只要求 $P$ 可逆。它允许进一步缩放变量，因此非零对角元可被归一化为 $1$ 或 $-1$，但对角元通常不再是原矩阵的特征值。

\begin{definition}[标准形与规范形]
二次型经过可逆变量替换得到的对角形
\[
d_1y_1^2+\cdots+d_ny_n^2
\]
通常称为标准形。若进一步把所有非零系数归一化为 $1$ 或 $-1$，得到
\[
y_1^2+\cdots+y_p^2-y_{p+1}^2-\cdots-y_{p+q}^2,
\]
则称为规范形。不同教材有时会交换这两个名称，因此判断题意时应以“是否要求系数归一化”为准。
\end{definition}

\begin{method}[两种方法的选择]
若题目要求“正交变换”“求主轴”“对角元为特征值”，应求特征值与标准正交特征向量。若只要求“化为标准形或规范形”“求惯性指数”“判断相合”，配方法通常计算更短。两种方法所得具体系数可能不同，但正、负、零项的个数相同。
\end{method}

\section{Sylvester 惯性定理与规范形}

\begin{theorem}[Sylvester 惯性定理]
对任意实对称矩阵 $A$，存在可逆矩阵 $P$ 使
\[
P^{\mathsf T}AP=\diag(I_p,-I_q,0_{n-p-q}).
\]
整数 $p,q$ 由 $A$ 唯一确定，分别称为正惯性指数和负惯性指数；$n-p-q$ 称为零惯性指数。并且
\[
p+q=\rank A.
\]
\end{theorem}
\begin{proof}
证明分为规范形的存在性与惯性指数的唯一性。

\textbf{存在性。} 相合对角化先给出可逆矩阵 $P_0$，使
\[
P_0^{\mathsf T}AP_0=D=\diag(d_1,\ldots,d_n).
\]
重新排列对角元后，设 $d_1,\ldots,d_p>0$，$d_{p+1},\ldots,d_{p+q}<0$，其余为零。取可逆对角矩阵
\[
S=\diag\left(\frac1{\sqrt{d_1}},\ldots,
\frac1{\sqrt{d_p}},
\frac1{\sqrt{|d_{p+1}|}},\ldots,
\frac1{\sqrt{|d_{p+q}|}},1,\ldots,1\right).
\]
逐个计算对角元可得
\[
S^{\mathsf T}DS=\diag(I_p,-I_q,0).
\]
因此取 $P=P_0S$，便有
\[
P^{\mathsf T}AP=S^{\mathsf T}(P_0^{\mathsf T}AP_0)S
=\diag(I_p,-I_q,0).
\]

\textbf{唯一性。} 设二次型在某组坐标下规范形为
\[
Q(y)=y_1^2+\cdots+y_p^2-y_{p+1}^2-\cdots-y_{p+q}^2.
\]
令 $r=p+q$。子空间
\[
V_+=\{(y_1,\ldots,y_p,0,\ldots,0)^{\mathsf T}\}
\]
维数为 $p$，且 $Q$ 在 $V_+\setminus\{0\}$ 上恒正。另一方面，若某子空间 $W$ 满足 $Q|_{W\setminus\{0\}}>0$，考虑线性映射
\[
\pi:W\to\R^p,\qquad \pi(y)=(y_1,\ldots,y_p)^{\mathsf T}.
\]
若 $\dim W>p$，由维数定理，线性映射 $\pi:W\to\R^p$ 不可能单射，因此 $\ker\pi$ 含非零向量 $w$。该向量的前 $p$ 个坐标全为零，从而
\[
Q(w)=-w_{p+1}^2-\cdots-w_{p+q}^2\le0,
\]
这与 $Q$ 在 $W\setminus\{0\}$ 上恒正矛盾。因此任何正定子空间的维数都不超过 $p$，而 $V_+$ 已达到维数 $p$。所以 $p$ 正是“$Q$ 为正定的子空间的最大维数”，它只由 $Q$ 本身决定，与所选坐标无关。

同理，后 $q$ 个负平方项对应的坐标子空间上 $Q$ 恒负；把任意负定子空间投影到这 $q$ 个坐标，重复上述核与维数的论证，可知 $q$ 是负定子空间的最大可能维数。因此 $p,q$ 唯一。

最后，相合变换保持秩，而规范形矩阵 $\diag(I_p,-I_q,0)$ 恰有 $p+q$ 个非零对角元，所以
\[
\rank A=p+q.
\]
\end{proof}

\begin{corollary}[实对称矩阵的相合分类]
两个同阶实对称矩阵相合，当且仅当它们的正、负、零惯性指数分别相同。
\end{corollary}
\begin{proof}
若 $A\cge B$，它们表示同一个二次型在两组可逆坐标下的矩阵，由惯性定理的唯一性，惯性指数相同。

反之，若 $A,B$ 的惯性指数同为 $(p,q,n-p-q)$，则存在可逆矩阵 $P,Q$ 使
\[
P^{\mathsf T}AP=J,\qquad Q^{\mathsf T}BQ=J,\qquad
J=\diag(I_p,-I_q,0).
\]
第二式等价于 $B=(Q^{-1})^{\mathsf T}JQ^{-1}$。代入第一式得
\[
B=(Q^{-1})^{\mathsf T}P^{\mathsf T}APQ^{-1}
=(PQ^{-1})^{\mathsf T}A(PQ^{-1}).
\]
$PQ^{-1}$ 可逆，故 $A\cge B$。
\end{proof}

\begin{corollary}[惯性指数与特征值符号]
实对称矩阵 $A$ 的正惯性指数、负惯性指数和零惯性指数，分别等于 $A$ 的正特征值、负特征值和零特征值的个数，均按代数重数计数。
\end{corollary}
\begin{proof}
由谱定理，存在正交矩阵 $U$ 使
\[
U^{\mathsf T}AU=\diag(\lambda_1,\ldots,\lambda_n).
\]
正交矩阵也是可逆矩阵，所以这是一个相合变换。将每个正的 $\lambda_i$、负的 $\lambda_i$ 分别缩放为 $1,-1$ 后得到惯性规范形。由惯性指数唯一性，三个惯性指数恰等于三类特征值的个数。
\end{proof}

\begin{corollary}[惯性指数对直和可加]
若实对称矩阵 $A,B$ 的正、负惯性指数分别为 $(p_A,q_A)$ 与 $(p_B,q_B)$，则分块对角矩阵
\[
A\oplus B=\begin{pmatrix}A&0\\0&B\end{pmatrix}
\]
的正、负惯性指数分别为
\[
(p_A+p_B,\ q_A+q_B).
\]
\end{corollary}
\begin{proof}
分别取可逆矩阵 $P,Q$，把 $A,B$ 化为各自的惯性规范形。则
\[
\diag(P,Q)^{\mathsf T}(A\oplus B)\diag(P,Q)
=(P^{\mathsf T}AP)\oplus(Q^{\mathsf T}BQ).
\]
右端的正平方项和负平方项个数分别相加，结论成立。
\end{proof}

\begin{corollary}[Witt 消去律]
设 $A_1,A_2,C$ 为实对称矩阵。若
\[
A_1\oplus C\cge A_2\oplus C,
\]
则 $A_1\cge A_2$。
\end{corollary}
\begin{proof}
相合矩阵具有相同的正、负、零惯性指数。由直和可加性，
\[
p(A_1)+p(C)=p(A_2)+p(C),\qquad
q(A_1)+q(C)=q(A_2)+q(C).
\]
消去两边相同的整数，得 $p(A_1)=p(A_2)$、$q(A_1)=q(A_2)$。矩阵阶数也相同，故零惯性指数相同。由实对称矩阵的相合分类，$A_1\cge A_2$。
\end{proof}

\begin{example}[title={利用配方法求惯性指数}]
设
\[
Q=x_1^2+4x_1x_2+2x_1x_3+2x_2^2+2x_2x_3+3x_3^2.
\]
求秩与惯性指数。
\begin{solution}
这就是前例的二次型，已化为
\[
Q=z_1^2-2z_2^2+\frac52z_3^2.
\]
故 $p=2$、$q=1$、零惯性指数为 $0$，从而 $\rank A=3$。进一步缩放可得规范形 $u_1^2+u_2^2-u_3^2$；平方项顺序不影响惯性指数。
\end{solution}
\end{example}

\section{二次型的符号与正定矩阵}

\begin{definition}[正定、半正定、负定与不定]
设 $A=A^{\mathsf T}\in\R^{n\times n}$，$Q(x)=x^{\mathsf T}Ax$。
\begin{enumerate}
\item 若对每个 $x\ne0$ 都有 $Q(x)>0$，则称 $A$ 正定，记作 $A>0$；
\item 若对每个 $x$ 都有 $Q(x)\ge0$，则称 $A$ 半正定，记作 $A\ge0$；
\item 若对每个 $x\ne0$ 都有 $Q(x)<0$，则称 $A$ 负定；
\item 若 $Q$ 既能取正值又能取负值，则称 $A$ 或 $Q$ 不定。
\end{enumerate}
\end{definition}

不能把“不是正定”直接理解为“不定”。例如 $\diag(1,0)$ 半正定但不正定，$\diag(-1,0)$ 半负定但不负定。

\begin{theorem}[惯性指数判据]
设 $A$ 为 $n$ 阶实对称矩阵，其惯性指数为 $(p,q,n-p-q)$。则
\begin{align*}
A>0 &\iff (p,q)=(n,0),\\
A\ge0 &\iff q=0,\\
A<0 &\iff (p,q)=(0,n),\\
A\text{ 不定} &\iff p>0\text{ 且 }q>0.
\end{align*}
\end{theorem}
\begin{proof}
作可逆变量替换后，二次型化为
\[
Q(y)=y_1^2+\cdots+y_p^2-y_{p+1}^2-\cdots-y_{p+q}^2.
\]
由于变量替换可逆，$x\ne0$ 当且仅当 $y\ne0$，所以符号性质可以直接在规范形上判断。

若 $p=n,q=0$，则 $Q(y)=\sum_{i=1}^n y_i^2$，对每个非零 $y$ 严格为正，故正定。反之，若 $q>0$，取仅第 $p+1$ 个坐标为 $1$ 的向量，得到 $Q=-1$；若存在零惯性方向，取对应单位向量得到 $Q=0$。这两种情况都与正定矛盾，因此正定必有 $(p,q)=(n,0)$。

对半正定性，若 $q=0$，规范形是若干平方和，故恒非负；若 $q>0$，沿任一负平方项对应的坐标轴取向量即可得到负值，所以不可能半正定。负定情形将 $Q$ 换成 $-Q$ 即得。不定要求同时取得正值和负值，这恰好等价于至少存在一个正平方项和一个负平方项，即 $p>0$ 且 $q>0$。
\end{proof}

\begin{theorem}[正定矩阵的常用等价条件]
设 $A=A^{\mathsf T}\in\R^{n\times n}$。下列条件等价：
\begin{enumerate}
\item $A$ 正定；
\item $A$ 的所有特征值均为正数；
\item $A$ 相合于 $I_n$；
\item 存在可逆实矩阵 $C$ 使 $A=C^{\mathsf T}C$；
\item $A$ 的所有顺序主子式 $\Delta_k=\det A[1{:}k,1{:}k]$ 均为正；
\item $A$ 的所有非空主子式均为正。
\end{enumerate}
\end{theorem}
\begin{proof}
证明按若干条短的等价链进行。

\textbf{(1)$\Longleftrightarrow$(2)。} 由实对称矩阵的谱定理，存在正交矩阵 $U$ 使
\[
U^{\mathsf T}AU=\diag(\lambda_1,\ldots,\lambda_n).
\]
令 $x=Uy$。由于 $U$ 可逆，$x\ne0$ 当且仅当 $y\ne0$，并且
\[
x^{\mathsf T}Ax=y^{\mathsf T}(U^{\mathsf T}AU)y
=\sum_{i=1}^n\lambda_i y_i^2.
\]
若所有 $\lambda_i>0$，该和对非零 $y$ 严格为正。反之，若某个 $\lambda_k\le0$，取 $y=e_k$，就有 $x^{\mathsf T}Ax=\lambda_k\le0$，与正定矛盾。

\textbf{(1)$\Longleftrightarrow$(3)。} 由惯性指数判据，$A$ 正定当且仅当其规范形为 $I_n$，这正是 $A\cge I_n$。

\textbf{(3)$\Longleftrightarrow$(4)。} 若 $P^{\mathsf T}AP=I$，在等式左乘 $(P^{-1})^{\mathsf T}$、右乘 $P^{-1}$，得到
\[
A=(P^{-1})^{\mathsf T}P^{-1}.
\]
取 $C=P^{-1}$ 即得 (4)。反之，若 $A=C^{\mathsf T}C$ 且 $C$ 可逆，则对 $x\ne0$，$Cx\ne0$，且
\[
x^{\mathsf T}Ax=\lVert Cx\rVert^2>0,
\]
故 (4) 推出 (1)。

\textbf{(1)$\Longrightarrow$(6)$\Longrightarrow$(5)。} 任取非空指标集 $I=\{i_1,\ldots,i_k\}$，记相应主子矩阵为 $A[I,I]$。对任意 $z\in\R^k\setminus\{0\}$，构造 $x\in\R^n$：在第 $i_1,\ldots,i_k$ 个位置依次放入 $z$ 的分量，其余位置为零。则
\[
z^{\mathsf T}A[I,I]z=x^{\mathsf T}Ax>0.
\]
故每个主子矩阵都正定，其全部特征值为正，行列式作为特征值之积也为正。因此所有非空主子式均正。顺序主子式只是主子式中的一部分，所以 (6) 推出 (5)。

\textbf{(5)$\Longrightarrow$(1)。} 使用无换行主元的对称消元。由于 $\Delta_1,\ldots,\Delta_n$ 均为正，特别地它们都非零，所以每一步消元的主元都非零，不需要交换行列。该消元得到唯一的分解
\[
A=LDL^{\mathsf T},\qquad d_k=\frac{\Delta_k}{\Delta_{k-1}},\quad \Delta_0=1.
\]
这里 $L$ 为单位下三角矩阵，$D=\diag(d_1,\ldots,d_n)$。公式 $d_k=\Delta_k/\Delta_{k-1}$ 可由左上 $k$ 阶分块的行列式得到：单位下三角矩阵的行列式为 $1$，故
\[
\Delta_k=d_1d_2\cdots d_k,
\]
再除以前一式 $\Delta_{k-1}=d_1\cdots d_{k-1}$ 即得。

因为所有 $\Delta_k>0$，所以每个 $d_k>0$。对 $x\ne0$，令 $y=L^{\mathsf T}x$。矩阵 $L^{\mathsf T}$ 可逆，故 $y\ne0$，并且
\[
x^{\mathsf T}Ax=y^{\mathsf T}Dy=\sum_{k=1}^n d_ky_k^2>0.
\]
故 $A$ 正定。
\end{proof}

\begin{corollary}[正定矩阵的 Cholesky 分解]
实对称矩阵 $A$ 正定，当且仅当存在对角元为正的可逆下三角矩阵 $R$，使
\[
A=RR^{\mathsf T}.
\]
在要求 $R$ 的对角元为正时，该分解唯一。
\end{corollary}
\begin{proof}
若 $A>0$，由上述 $LDL^{\mathsf T}$ 分解可知 $D$ 的对角元全为正。令
\[
D^{1/2}=\diag(\sqrt{d_1},\ldots,\sqrt{d_n}),\qquad R=LD^{1/2}.
\]
则
\[
RR^{\mathsf T}=LD^{1/2}D^{1/2}L^{\mathsf T}=LDL^{\mathsf T}=A.
\]
$R$ 是下三角矩阵且对角元为正。反之，若 $A=RR^{\mathsf T}$ 且 $R$ 可逆，则对 $x\ne0$，
\[
x^{\mathsf T}Ax=\lVert R^{\mathsf T}x\rVert^2>0.
\]

再证唯一性。若另有对角元为正的可逆下三角矩阵 $S$ 满足
\[
A=SS^{\mathsf T}=RR^{\mathsf T},
\]
令 $Q=R^{-1}S$。下三角矩阵的逆及乘积仍为下三角矩阵，所以 $Q$ 为下三角矩阵；并且
\[
QQ^{\mathsf T}
=R^{-1}SS^{\mathsf T}(R^{-1})^{\mathsf T}
=R^{-1}RR^{\mathsf T}(R^{-1})^{\mathsf T}=I.
\]
故 $Q$ 是正交矩阵。下三角正交矩阵的各列两两正交且长度为 $1$：从最后一列开始，该列只有最后一个分量可能非零，故其绝对值为 $1$；再利用它与前面各列正交，可依次推出最后一行的其余元素为零。逐阶重复，得到 $Q$ 必为对角矩阵，且每个对角元为 $\pm1$。由于 $R,S$ 的对角元均为正，$Q=R^{-1}S$ 的对角元也为正，因此 $Q=I$，从而 $R=S$。
\end{proof}

\begin{remark}
顺序主子式仅有 $n$ 个：左上角 $1,2,\ldots,n$ 阶主子式。主子式则由任意相同的行、列指标集选出。Sylvester 判据表明，检验实对称矩阵正定只需检查顺序主子式；“所有主子式均正”虽等价，但计算量更大。
\end{remark}

\begin{corollary}[负定的 Sylvester 判据]
设 $A=A^{\mathsf T}$，$\Delta_k$ 为其 $k$ 阶顺序主子式。则
\[
A<0
\quad\Longleftrightarrow\quad
(-1)^k\Delta_k>0\quad(k=1,\ldots,n).
\]
等价地，$\Delta_1<0,\Delta_2>0,\Delta_3<0,\ldots$，符号正负交替。
\end{corollary}
\begin{proof}
$A<0$ 当且仅当 $-A>0$。矩阵 $-A$ 的左上 $k$ 阶主子矩阵为 $-A_k$，其行列式为
\[
\det(-A_k)=(-1)^k\det A_k=(-1)^k\Delta_k.
\]
对 $-A$ 使用正定的 Sylvester 判据即可。
\end{proof}

\begin{theorem}[半正定与主子式判据]
设 $A=A^{\mathsf T}\in\R^{n\times n}$。则
\[
A\ge0
\quad\Longleftrightarrow\quad
A\text{ 的所有主子式均非负}.
\]
类似地，$A\le0$ 当且仅当每个 $k$ 阶主子式 $M_k$ 都满足 $(-1)^kM_k\ge0$。
\end{theorem}
\begin{proof}
先设 $A\ge0$。任取指标集 $I$，主子矩阵 $A[I,I]$ 的二次型是把其余坐标置零后得到的限制，所以仍半正定。实对称半正定矩阵的特征值均非负，因此 $\det A[I,I]$ 作为这些特征值的乘积也非负。

反之，设 $A$ 的所有主子式均非负。考虑多项式
\[
p(t)=\det(tI_n+A).
\]
$t^{n-k}$ 的系数等于 $A$ 的所有 $k$ 阶主子式之和，因此 $p(t)$ 的所有系数均非负，且最高次项系数为 $1$。于是对每个 $t>0$ 都有 $p(t)>0$。若 $A$ 存在负特征值 $\lambda<0$，取 $t=-\lambda>0$，则 $tI_n+A$ 有零特征值，从而 $p(t)=0$，矛盾。故 $A$ 没有负特征值，即 $A\ge0$。

半负定结论对 $-A$ 使用上一结论即可，因为 $-A$ 的每个 $k$ 阶主子式等于原主子式乘以 $(-1)^k$。
\end{proof}

\begin{pitfall}[半正定判据必须检查所有主子式]
正定矩阵只需检查顺序主子式严格为正；半正定时，不能只把条件机械改成“顺序主子式均非负”。例如
\[
\begin{pmatrix}0&0\\0&-1\end{pmatrix}
\]
的顺序主子式为 $0,0$，均非负，但矩阵并不半正定。失败原因是遗漏了指标集 $\{2\}$ 对应的主子式 $-1$。
\end{pitfall}

\begin{theorem}[半正定矩阵的平方根与分解]
设 $A=A^{\mathsf T}\in\R^{n\times n}$。下列条件等价：
\begin{enumerate}
\item $A\ge0$；
\item 存在唯一的实对称半正定矩阵 $S$，使 $S^2=A$；
\item 存在实矩阵 $C$，使 $A=C^{\mathsf T}C$。
\end{enumerate}
矩阵 $S$ 称为 $A$ 的主平方根，记作 $A^{1/2}$。此外，$A>0$ 当且仅当 $A^{1/2}$ 可逆。
\end{theorem}
\begin{proof}
\textbf{(1)$\Longrightarrow$(2)。} 由谱定理，存在正交矩阵 $U$ 使
\[
A=U\diag(\lambda_1,\ldots,\lambda_n)U^{\mathsf T},
\qquad \lambda_i\ge0.
\]
定义
\[
S=U\diag(\sqrt{\lambda_1},\ldots,\sqrt{\lambda_n})U^{\mathsf T}.
\]
则 $S$ 对称、半正定，且 $S^2=A$。

若另有 $T\ge0$ 且 $T^2=A$，则 $T$ 与 $A=T^2$ 交换，所以 $T$ 保持 $A$ 的每个特征子空间。在特征值 $\lambda$ 对应的子空间上，$T^2=\lambda I$；由于 $T$ 半正定，$T$ 的特征值只能为 $\sqrt\lambda$。故 $T$ 在该子空间上等于 $\sqrt\lambda I$，从而 $T=S$。

\textbf{(2)$\Longrightarrow$(3)。} 取 $C=S$，则 $A=S^2=S^{\mathsf T}S$。

\textbf{(3)$\Longrightarrow$(1)。} 对任意 $x$，
\[
x^{\mathsf T}Ax=\lVert Cx\rVert^2\ge0.
\]
最后，$A$ 正定当且仅当所有 $\lambda_i>0$，这等价于 $A^{1/2}$ 可逆。
\end{proof}

\begin{theorem}[Rayleigh 商与特征值界]
设 $A=A^{\mathsf T}$ 的最小、最大特征值分别为 $\lambda_{\min}$、$\lambda_{\max}$。对任意 $x\ne0$，定义
\[
R_A(x)=\frac{x^{\mathsf T}Ax}{x^{\mathsf T}x}.
\]
则
\[
\lambda_{\min}\le R_A(x)\le\lambda_{\max}.
\]
其最小值、最大值分别等于 $\lambda_{\min},\lambda_{\max}$，并在相应特征子空间的非零向量处取得。
\end{theorem}
\begin{proof}
取正交矩阵 $U$ 使 $U^{\mathsf T}AU=\diag(\lambda_1,\ldots,\lambda_n)$，令 $x=Uy$。则
\[
R_A(x)
=\sum_{i=1}^n\lambda_i
\frac{y_i^2}{\sum_{j=1}^n y_j^2}.
\]
各权重非负且总和为 $1$，所以 Rayleigh 商是全部特征值的凸组合。取相应特征向量即可达到两个端点。
\end{proof}

\begin{corollary}[正定性的统一下界]
实对称矩阵 $A$ 正定，当且仅当存在 $c>0$，使
\[
x^{\mathsf T}Ax\ge c\lVert x\rVert^2
\]
对所有 $x$ 成立。最大的可行常数为 $\lambda_{\min}(A)$。
\end{corollary}
\begin{proof}
若 $A>0$，则 $\lambda_{\min}>0$，Rayleigh 商定理给出结论。反之，该不等式对每个非零 $x$ 给出 $x^{\mathsf T}Ax>0$。取最小特征值对应的单位特征向量，可知任何可行的 $c$ 都不超过 $\lambda_{\min}$。
\end{proof}

\begin{theorem}[Schur 补正定判据]
设
\[
A=\begin{pmatrix}B&C\\C^{\mathsf T}&D\end{pmatrix}
\]
为实对称分块矩阵。
\begin{enumerate}
\item 若 $B$ 可逆，则
\[
A>0
\quad\Longleftrightarrow\quad
B>0\ \text{且}\ D-C^{\mathsf T}B^{-1}C>0.
\]
\item 若 $D$ 可逆，则
\[
A>0
\quad\Longleftrightarrow\quad
D>0\ \text{且}\ B-CD^{-1}C^{\mathsf T}>0.
\]
\end{enumerate}
两个差矩阵分别称为关于 $B$ 和 $D$ 的 Schur 补。
\end{theorem}
\begin{proof}
只证第一条，第二条交换分块顺序即可。取可逆分块矩阵
\[
P=\begin{pmatrix}I&-B^{-1}C\\0&I\end{pmatrix}.
\]
利用 $B^{\mathsf T}=B$ 可得 $(B^{-1})^{\mathsf T}=B^{-1}$。直接相乘得到
\[
P^{\mathsf T}AP=
\begin{pmatrix}
B&0\\
0&D-C^{\mathsf T}B^{-1}C
\end{pmatrix}.
\]
相合变换保持正定性：若 $x=Py$ 且 $P$ 可逆，则非零 $x$ 与非零 $y$ 一一对应，并有
\[
y^{\mathsf T}(P^{\mathsf T}AP)y=x^{\mathsf T}Ax.
\]
分块对角矩阵正定，当且仅当两个对角块都正定，因此结论成立。
\end{proof}

\begin{theorem}[Jacobi 惯性判据]
设 $A=A^{\mathsf T}$ 的所有顺序主子式
\[
\Delta_1,\Delta_2,\ldots,\Delta_n
\]
均非零，并约定 $\Delta_0=1$。则存在 $LDL^{\mathsf T}$ 分解，其中
\[
D=\diag\left(
\frac{\Delta_1}{\Delta_0},
\frac{\Delta_2}{\Delta_1},
\ldots,
\frac{\Delta_n}{\Delta_{n-1}}
\right).
\]
因此，$A$ 的正、负惯性指数分别等于数列
\[
1,\Delta_1,\Delta_2,\ldots,\Delta_n
\]
中相邻两项同号和异号的次数。
\end{theorem}
\begin{proof}
顺序主子式均非零保证对称消元过程中每个主元都非零，因此无需交换行列，可得
\[
A=LDL^{\mathsf T},
\]
其中 $L$ 为单位下三角矩阵，$D=\diag(d_1,\ldots,d_n)$。比较左上 $k$ 阶分块的行列式：
\[
\Delta_k=d_1d_2\cdots d_k,
\]
所以
\[
d_k=\frac{\Delta_k}{\Delta_{k-1}}.
\]
矩阵 $A$ 与 $D$ 相合，故惯性指数等于 $D$ 中正、负对角元的个数。若 $\Delta_{k-1}$ 与 $\Delta_k$ 同号，则 $d_k>0$；若异号，则 $d_k<0$，结论成立。
\end{proof}

\begin{remark}
Jacobi 惯性判据要求所有顺序主子式非零。若中间某个 $\Delta_k=0$，商 $\Delta_{k+1}/\Delta_k$ 无定义，不能直接数符号变化；此时应改用带主元交换的对称消元、配方法或特征值方法。
\end{remark}

\begin{example}[title={含参数矩阵的正定范围}]
求实数 $a$ 的范围，使
\[
A=\begin{pmatrix}
2&1&0\\
1&a&1\\
0&1&2
\end{pmatrix}
\]
正定。
\begin{solution}
由 Sylvester 判据，
\[
\Delta_1=2>0,\qquad \Delta_2=2a-1,\qquad \Delta_3=4a-4.
\]
故 $A>0$ 当且仅当 $2a-1>0$ 且 $4a-4>0$，即 $a>1$。第二个条件已经蕴含第一个条件，但书写时仍应列出全部顺序主子式。
\end{solution}
\end{example}

\begin{example}[title={判断半正定不能只看顺序主子式严格不等式}]
判断
\[
A=\begin{pmatrix}0&0\\0&1\end{pmatrix}
\]
的符号性质。
\begin{solution}
$x^{\mathsf T}Ax=x_2^2\ge0$，故 $A$ 半正定，但取 $x=(1,0)^{\mathsf T}\ne0$ 时二次型为零，所以不正定。半正定的一个判据是所有主子式均非负；只把 Sylvester 正定判据中的“$>$”机械改成“$\ge$”并只检查顺序主子式，并不足以判定半正定。
\end{solution}
\end{example}

\section{对称双线性型的对角化与规范形}

设 $B$ 是有限维实线性空间 $V$ 上的对称双线性型。任取一组基 $\mathcal E$，其矩阵 $G$ 对称。由相合对角化，存在可逆矩阵 $P$ 使 $P^{\mathsf T}GP$ 为对角阵。取新基 $\widetilde{\mathcal E}=\mathcal E P$，便得到
\[
B(\widetilde e_i,\widetilde e_j)=0\qquad(i\ne j).
\]
这种基称为 $B$ 的正交基；这里“正交”是相对于 $B$ 而言，不一定是欧氏内积下的正交。

\begin{definition}[根空间与非退化性]
对称双线性型 $B$ 的根空间定义为
\[
\operatorname{rad}B
=\{u\in V:B(u,v)=0\text{ 对所有 }v\in V\}.
\]
若 $\operatorname{rad}B=\{0\}$，则称 $B$ 非退化；否则称 $B$ 退化。
\end{definition}

\begin{proposition}[根空间、矩阵核与秩]
若 $G$ 是 $B$ 在任一组基下的 Gram 矩阵，则
\[
[\operatorname{rad}B]_{\mathcal E}=\ker G,\qquad
\dim\operatorname{rad}B=n-\rank G.
\]
因此可定义 $\rank B=\rank G$，该数与基的选取无关。
\end{proposition}
\begin{proof}
设 $x=[u]_{\mathcal E}$。向量 $u$ 属于 $\operatorname{rad}B$，当且仅当对每个坐标向量 $y$ 都有
\[
x^{\mathsf T}Gy=0.
\]
这等价于行向量 $x^{\mathsf T}G=0$。因为 $G$ 对称，这又等价于 $Gx=0$，即 $x\in\ker G$。维数公式由秩零度定理得到。换基后 $\widetilde G=P^{\mathsf T}GP$，左右乘可逆矩阵不改变秩，所以 $\rank B$ 与基无关。
\end{proof}

\begin{definition}[$B$-正交补]
设 $W$ 是 $V$ 的子空间。定义
\[
W^{\perp_B}
=\{v\in V:B(v,w)=0\text{ 对所有 }w\in W\}.
\]
对一般对称双线性型，$W$ 可能与 $W^{\perp_B}$ 有非零交集。
\end{definition}

\begin{theorem}[非退化双线性型的正交补]
若 $B$ 在有限维空间 $V$ 上非退化，则对任意子空间 $W$，
\[
\dim W+\dim W^{\perp_B}=\dim V,
\qquad
(W^{\perp_B})^{\perp_B}=W.
\]
若 $B$ 在 $W$ 上的限制也非退化，则
\[
V=W\oplus W^{\perp_B}.
\]
\end{theorem}
\begin{proof}
定义
\[
\Phi:V\to W^*,\qquad \Phi(v)(w)=B(v,w).
\]
其核恰为 $W^{\perp_B}$。因为 $B$ 非退化，映射 $T_B:V\to V^*$，$T_B(v)=B(v,\cdot)$ 是同构。任意 $W$ 上的线性泛函都可延拓到 $V$，再由 $T_B$ 的满射性写成 $B(v,\cdot)$，所以 $\Phi$ 满射。由秩零度定理，
\[
\dim V=\dim W^{\perp_B}+\dim W.
\]
显然 $W\subseteq(W^{\perp_B})^{\perp_B}$，而两者维数相同，故相等。

若 $B|_W$ 非退化，则 $W\cap W^{\perp_B}=\{0\}$。再结合维数之和为 $\dim V$，得到直和分解。
\end{proof}

\begin{corollary}[退化部分与非退化部分的分解]
对任意有限维实对称双线性型 $B$，都存在子空间 $W$，使
\[
V=\operatorname{rad}B\oplus W,
\]
且 $B|_W$ 非退化。
\end{corollary}
\begin{proof}
在规范基下，令 $W$ 由对应于 $1$ 与 $-1$ 的基向量张成，根空间由对应于零项的基向量张成。两部分基向量合起来构成 $V$ 的一组基；$B|_W$ 的 Gram 矩阵为 $\diag(I_p,-I_q)$，可逆，故非退化。
\end{proof}

\begin{theorem}[对称双线性型的规范形]
存在 $V$ 的一组基，使
\[
B(u,v)=\sum_{i=1}^p u_iv_i-\sum_{i=p+1}^{p+q}u_iv_i,
\]
其余坐标不出现。整数 $p,q$ 与基的选取无关。
\end{theorem}
\begin{proof}
任取 $V$ 的一组基 $\mathcal E$，设 $B$ 在该基下的 Gram 矩阵为 $G$。由于 $B$ 对称，$G$ 是实对称矩阵。由 Sylvester 惯性定理，存在可逆矩阵 $P$ 使
\[
P^{\mathsf T}GP=J=\diag(I_p,-I_q,0).
\]
定义新基 $\widetilde{\mathcal E}=\mathcal E P$。$P$ 可逆保证新向量组仍是基。由基变换公式，$B$ 在新基下的矩阵正是 $J$。

若
\[
[u]_{\widetilde{\mathcal E}}=(u_1,\ldots,u_n)^{\mathsf T},\qquad
[v]_{\widetilde{\mathcal E}}=(v_1,\ldots,v_n)^{\mathsf T},
\]
则
\begin{align*}
B(u,v)
&=[u]_{\widetilde{\mathcal E}}^{\mathsf T}
J[v]_{\widetilde{\mathcal E}}\\
&=u_1v_1+\cdots+u_pv_p
-u_{p+1}v_{p+1}-\cdots-u_{p+q}v_{p+q}.
\end{align*}
惯性定理还表明 $p,q$ 与最初选取的基无关，因此它们是双线性型 $B$ 本身的不变量。
\end{proof}

\begin{corollary}[退化性与零惯性指数]
对称双线性型 $B$ 非退化，当且仅当其规范形中没有零项，即 $p+q=n$。
\end{corollary}
\begin{proof}
在规范基下，Gram 矩阵为 $J=\diag(I_p,-I_q,0_{n-p-q})$。双线性型非退化等价于 Gram 矩阵可逆，而 $J$ 可逆当且仅当零块大小为 $0$，即 $p+q=n$。
\end{proof}

若 $q=0$ 且 $p=n$，则 $B$ 是内积；此时可进一步利用 Gram--Schmidt 过程选取一组 $B$-标准正交基，使 Gram 矩阵为 $I_n$。若 $B$ 不正定，传统的标准正交化中会出现零长度或负平方，不能不加修改地套用欧氏内积版本。

\begin{example}[title={三个典型双线性型}]
下列映射均为双线性型：
\begin{enumerate}
\item $\R^2$ 上 $B(x,y)=2x_1y_1+3x_1y_2+3x_2y_1+5x_2y_2$，其矩阵为 $\begin{pmatrix}2&3\\3&5\end{pmatrix}$；
\item 矩阵空间 $\R^{m\times n}$ 上 $B(X,Y)=\tr(X^{\mathsf T}Y)$，它实际上是 Frobenius 内积；
\item 连续函数空间 $C[a,b]$ 上 $B(f,g)=\int_a^b f(t)g(t)\,\mathrm dt$。该空间无限维，但双线性型的定义仍成立。
\end{enumerate}
\end{example}

\begin{summary}[本章知识链]
二次齐次多项式通过“交叉项系数折半”得到唯一对称矩阵；对称矩阵给出二次型 $Q(x)=x^{\mathsf T}Ax$ 和对称双线性型 $B(x,y)=x^{\mathsf T}Ay$；极化公式在二者之间往返。可逆变量替换或基变换使矩阵按 $P^{\mathsf T}AP$ 变化，即相合变换。配方法给出相合对角化，Sylvester 惯性定理把对角系数进一步归一化为 $1,-1,0$，正定性则等价于规范形中只有 $n$ 个正平方项。
\end{summary}
